% Compression rewrite (DECISION-20261119c; REWRITE-RULES-20261119.md rule 5): text the limitation-resolver lane
% moved verbatim out of the posted note, with one move-manifest row per sentence (LR-C ids in
% history-cuts/move-manifest-lr-20261117.tsv). References into the note resolve through main_note.aux (xr).
\section{Method and statistics detail moved from the note}
This section holds the procedure and supporting detail that the compressed note no longer prints, as the note printed it.

\subsection{Truth-level W residual (from Part 2, the five observables)}
% From sec_observables.tex:40 (draft f6a4fc372 lineage); detail a reader does not need to interpret a result
At truth level, $q_{3}$ is computed from the true neutrino and muon four-vectors, but $\theta_{\mu}$, and with it $\pT$ and $\pPar$, is measured from a fixed beam axis tilted by \SI{58.87}{mrad} from the detector axis rather than from the true neutrino direction.  Truth $W$ is therefore fixed by $(\pT,\pPar,q_{3})$ only up to that event-by-event angle, of about \SI{1}{mrad}, and the choice of $M$: the residual is within \SI{\pm2}{MeV} for 68\% of events and within $-49$ to $+67$~MeV for 99.9\%; \SI{0.25}{\percent} of events (weighted) change their coarse $W$ bin, and \SI{0.086}{\percent} (weighted) have no real solution. % D1 (crossdoc audit 2026-11-15): weighted, probe_omni.json:10 q3sq_minus_pt2_negative_wfrac 0.0008646 (the 0.15% was the unweighted 49754/32846302, :9 and :8); the Supplement prints m4w-noreal-pct, 0.086
% Methods rewrite (M4, the truth-W residual): M4-W-RECORD.md :6 (theta_true from a fixed beam axis tilted -58.87 mrad
% about x, MAT TruthFunctions.h:21-27, PlotUtilsPhysicalConstants.h:44), :7 and :74-81 (W_true - W_pred with M_bar: 68%
% within about +-2 MeV, 99.9% within -49/+67 MeV), :85-90 (0.25% of events, weighted, change coarse W bin; 0.15% have
% no real solution), :63 (the tilt is the only other cause), :55-59 (median 1.03 mrad between the true neutrino and that
% axis); q3 from the true four-vectors, README-gen5d-build.md:50-58. REVIEW-1 finding 18: M4-W-RECORD.md:7, :72.

\subsection{Per-event analysis trees (from Part 2, sec:evloop)}
% From sec_experiment.tex:34 (draft f6a4fc372 lineage); procedure a reader does not need to interpret a result
Unbinned unfolding needs per-event records rather than histograms. A dedicated event loop (App.~\ref{app:codebase}), built on MAT's systematic-universe tools, writes four per-event samples for each playlist: all true signal events (the efficiency denominator), reconstructed truth-matched simulated events, simulated background, and data. Entries carry truth and reconstructed observables ($\pT$, $\pPar$, and in extended runs $\Eavail$, $q_{3}$, $W$) and the MINERvA Tune~v1 central-value weight. The systematic-universe samples also carry every universe's weight and the shifted observables of each lateral band, a variation that shifts reconstructed kinematics rather than reweighting events (a vertical band; \S\ref{sec:systematics}). Two conditions make the truth sample the complete reference for the simulated signal:
% (\texttt{\_universes\_full})
\begin{itemize}
  \item A reconstructed signal event is filled only if its $(\mathrm{run},\mathrm{subrun},\mathrm{event})$ key is in the truth denominator. Truth events without a reconstructed partner enter as native OmniFold ``miss'' entries.
  \item The simulated signal sample is complete by construction at ME~FHC (completeness $c=1.000000$; \S\ref{sec:method-xsec}).
\end{itemize}
The loop reads one playlist-wide flux and calibration from the first event, so it is run separately for each playlist, and the twelve per-playlist outputs are combined. The corresponding baseline histogram event loop supplies the flux histogram, POT-weighted across playlists; the size of the playlist-to-playlist flux differences is stated in \S\ref{sec:limitations-experiment}.
% (\texttt{TTree::SetMaxTreeSize}-aware)

\subsection{Shapes of the reconstructed control distributions (from Part 2, sec:controlplots)}
% From sec_experiment.tex:49 (draft f6a4fc372 lineage); description of the control figures beyond their captions
The beam and acceptance set the shapes. The $\pPar$ peak at $4$--$6$~GeV/$c$ reflects the ${\langle E_\nu\rangle}\approx\SI{6}{GeV}$ NuMI medium-energy peak; spectrometer matching removes the lowest-$\pPar$ muons. The $\pT$ spectrum peaks near $0.7$~GeV/$c$ and falls by over an order of magnitude across the measured range. Low-recoil (quasi-elastic and $2p2h$) events dominate the sample. Available energy $\Eavail$ concentrates below ${\sim}\SI{0.5}{GeV}$, with a long inelastic tail in the wide upper bins; $q_{3}$ tracks this population's momentum transfer. Reconstructed $W$ rises through the smeared quasi-elastic and $\Delta(1232)$ region before the deep-inelastic continuum takes over above ${\sim}\SI{2}{GeV}$. Absolute background is at the few-percent level under the muon projections and concentrates at high recoil (highest $\Eavail$, $q_{3}$ and $W$ bins).

\subsection{Iteration count and the Huang et al.\ recommendation (from Part 2, sec:method-config)}
% From sec_method.tex (dfc41cdf5) (draft f6a4fc372 lineage); rationale a reader does not need to interpret a result
Huang \emph{et al.}~\cite{Huang:2025omnifold} select 45 iterations for OmniFold-family methods (20 for their binned benchmark) using a binned $\chi^{2}$ against known mock-data truth, a criterion they note is infeasible on real data. We do not adopt that recommendation. Their mock data injects a large prior--data \emph{shape} discrepancy through an RPA-modified quasielastic model, requiring many iterations. The simulation prior differs from the measured distributions in shape as well as normalization: the data exceed it by about 12\%, and the data/MC ratio rises with $\pT$ from 0.88 to about 1.2 (\S\ref{sec:controlplots}). Cross-experiment production also typically uses ${\sim}5$ iterations~\cite{Canelli:2025guide}.

\subsection{MINERvA driver and pipeline conventions (from Part 2, sec:method-driver)}
% From sec_method.tex (dfc41cdf5) (draft f6a4fc372 lineage); procedure a reader does not need to interpret a result
The higher-dimensional MINERvA driver (App.~\ref{app:codebase}) extends the 2D driver using an axis list and the same 2D helpers. The first two axes are always $(\pT,\pPar)$, carrying the fiducial phase-space cut and the per-$\pT$ flux. Higher-dimensional runs extend the muon observables $(\pT,\pPar)$ with $\Eavail$, $q_{3}$, $W$ in the phase space of \S\ref{sec:selection}. Every production, universe-sweep, and closure run follows these conventions:
\begin{itemize}[itemsep=2pt]
  \item \textbf{Reco-space masking.} The measured events entering step-1 training are exactly those passing the reconstructed-level selection; nothing outside the reconstructed phase space reaches the classifier.
  \item \textbf{Step-1 training within the phase space.} Step 1 targets the background-subtracted data of \S\ref{sec:method-algo}, purity-weighted or with background simulation injected at negative weight (refined to non-negative weights in the refined configurations). Genuine backgrounds \emph{and} out-of-phase-space signal are removed from the data side. The simulated reconstructed sample contains no out-of-phase-space signal by construction, so neither side of the step-1 training contains events whose truth lies outside the measurement phase space.
        % Accuracy fix (unscored; director ruling 2026-10-26): the list covers every production run, and the 2D result
        % and the 3D and 4D central values use refined negative-weight subtraction (sec:method-algo). Mechanism:
        % 2d-unfolding/unfold_2d_omnifold_unbinned.py:1362-1377 and :1825-1870 @ ddf16f10; nd-unfolding/unfold_nd_omnifold_unbinned.py:619-630 and :900-1030 @ ddf16f10.
        % In the negative-weight modes the signal fakes are injected with the background (2d-unfolding/unfold_2d_omnifold_unbinned.py:1846-1858 @ ddf16f10).
  \item At histogramming, each simulated truth event carries its final OmniFold push-weight times its central-value simulation weight, including MINERvA Tune~v1 reweighting. No efficiency-histogram division is applied: the step-1 regression for unreconstructed truth events already gives an efficiency-corrected truth spectrum, and dividing again would double-correct.
        % Development history: a failure mode that occurred once in
        % development and was caught by the marginal self-validation gate
        % (\S\ref{sec:validation}).
  \item \textbf{Closure mirrors production.} In closure tests the pseudo-data are simulated events passing both the reconstructed and truth selections, weighted identically to the production data path, so closure exercises the same code path as the data unfolding.
\end{itemize}

\subsection{Structure of the truth-to-reco migrations (from Part 2, sec:method-binning)}
% From sec_method.tex (dfc41cdf5) (draft f6a4fc372 lineage); description of the migration figure beyond its caption
The structure differs qualitatively between the muon and hadron sides. The spectrometer muon measurement makes the $\pT$ and $\pPar$ maps nearly tridiagonal, with migrations rarely exceeding one bin. Reconstruction stability dips to ${\sim}0.45$--$0.5$ for $\pT$ and ${\sim}0.35$ for $\pPar$ where bins are finest, recovering in wide edge bins. The calorimetric axes have coarser bins but relatively broader, visibly asymmetric migrations. $\Eavail$ preferentially migrates \emph{downward} because energy carried by neutrons and other unresponsive particles is not collected. Reconstruction stability falls from $0.61$ in the lowest bins to $0.43$ in the highest. $q_{3}$, reconstructed from the muon $Q^{2}$ and the calorimetric recoil, inherits a milder version of the same pattern. $W$ compounds $\Eavail$ and muon smearing and has the broadest migrations; two-bin shifts are common in the resonance region. The lowest-$W$ (quasi-elastic) bin has the highest stability at $0.71$.

\subsection{No underflow or overflow in the reported product (from Part 2, sec:method-binning)}
% From sec_method.tex @1ad08b466; argument a reader does not need to interpret a result
Neither underflow nor overflow appears in the reported product, by construction rather than by truncation: the unbinned unfolding has no response-matrix under- or overflow, and the reporting bins are applied once, to events that all lie inside the grids. Truth-level muon edges coincide with the measurement phase space (\S\ref{sec:selection}); no signal is defined beyond them. Hadronic axes span $[0,100]\,\si{GeV}$, with catch-bin edges far beyond the kinematic reach at $\langle E_\nu\rangle\sim\SI{6}{GeV}$, so the truth grids are exhaustive. A selected reco event outside the grid is treated as unaccepted and excluded from the measured sample by \S\ref{sec:method-driver}'s masking convention. Its in-phase-space truth content is recovered, as for any unreconstructed event, by the treatment of truth events without a reconstructed partner.

\subsection{Profile OmniFold~\cite{Zhu:2025profile} and related developments (from Part 2)}
% From sec_method.tex @1ad08b466; literature context
Its distinct algorithmic construction complements earlier unbinned profiled unfolding~\cite{Chan:2023tbf}.

\subsection{Gradient-boosted trees in unbinned unfolding (from Part 2, sec:method-config)}
% From sec_method.tex @1ad08b466; literature context
Ref.~\cite{Milton:2025tools} demonstrates their effectiveness for unbinned unfolding with few high-level features, complementing deep models for low-level event representations.

\subsection{MINERvA's published-covariance description (from App. A, sec:minerva)}
% From app_statmethods.tex:235 @143f4fc97; context a reader does not need to interpret a result
% Framing added in the records (not moved text).
\noindent\emph{(``Below'' points to the note's MAT paragraph of \S\ref{sec:minerva}, which states the CV-centered pair formula.)}
\paragraph{Ruterbories 2021, Section V.} The paper evaluates three uncertainty categories (flux, detector response and the neutrino-interaction model) by ``re-extracting the cross section using modified simulations'', i.e.\ the universe method, and gives no covariance formula. It reports a flux uncertainty below $4\%$ across most of phase space as a result of the $\nu$-$e$-scattering constraint (Valencia, arXiv:1906.00111); a $1.4\%$ target-nucleon normalization; an in-situ test-beam hadron response with per-particle GEANT inelastic variations; and a $1\%$ muon energy scale constrained by the low-recoil fit of Bashyal (JINST 16 P08068). It states that \emph{``the fit causes the flux and muon energy scale uncertainties to be correlated and that correlation is propagated to the final result''}; this is the only correlation between bands that the paper specifies. The only covariance in its equations is $V$ in the model-test $\chi^{2}$ (Eq.~1, ``the measurement covariance matrix''), with $205$ degrees of freedom (Table I); the paper also gives a log-normal $\chi^{2}$ variant, motivated by the highly correlated flux normalization (Peelle's Pertinent Puzzle). The paper itself gives no per-band covariance construction, no $\pm1\sigma$ pair formula and no band-independence or block-sum statement. MAT's code supplies the per-band construction and the independent band sum, but no separate pair formula (below). Nor does it state the rank, condition number and regularization of $V$, or any bootstrap or coverage validation (it uses D'Agostini iterative unfolding).

\subsection{Upstream references of the band amplitudes (from App. A, sec:minerva)}
% From app_statmethods.tex:243 @143f4fc97; reference list
\paragraph{Upstream references that set each band's $\pm1\sigma$ amplitude.}
\begin{itemize}[nosep]
  \item Flux / PPFX (100 universes), Aliaga, PRD 94 092005 (2016), arXiv:1607.00704;
  \item $\nu$-$e$ flux constraint ($<4\%$), Valencia, PRD 100 092001 (2019), arXiv:1906.00111;
  \item muon-E scale $\leftrightarrow$ flux fit, Bashyal, JINST 16 (2021) P08068;
  \item test-beam / GEANT hadron response, Aliaga, NIM A 789 28 (2015), arXiv:1501.06431;
  \item GENIE reweighting parameters, Andreopoulos, NIM A 614 87 (2010), and the GENIE manual arXiv:1510.05494;
  \item $1.4\%$ normalization (detector mass), Aliaga, NIM A 743 130 (2014), arXiv:1305.5199.
\end{itemize}

\subsection{Log-normal chi2 definitions compared on an earlier covariance (from App. A, sec:chisq)}
% From app_statmethods.tex:335 @143f4fc97; history of an earlier covariance
% Framing added in the records (not moved text); the J3/J4 comment that the moved "comment below" names stays in app_statmethods.tex.
\noindent\emph{(``The definition above'' is the note's log-normal definition in \S\ref{sec:chisq}.)}
 An independent recomputation compared five definitions on an earlier covariance, before the refined estimator was adopted. Only the definition above reproduced the value then quoted. The other four gave $1.197$ (our covariance scaled by our own values), $0.958$ (both scaled by ours), $0.963$ (scaled by the geometric mean) and $1.015$ ($\ln(1+\mathbf{C}/(x_{i}x_{j}))$ in place of the ratio). The quoted $\chiCombined$ and $\chiCombinedLog$ have no independent recomputation: they are the production run's own outputs. Its checks confirm that the covariance's reported cells are the paper's and that the plotted projections equal the budget's, and neither recomputes these values. Adding the OmniFold error budget lowers the summed-covariance distance from $\chiPaper$ to $\chiCombined$ per bin. % G2R no-recompute sentence (the director, 2026-11-15): RE-2-shapes run 2 has no independent chi2 recompute (cov2d, 2026-11-15); its check_rebuild (g2r/logs/check_rebuild.log, 5/5) checks the chi2 files' cell identity (chi2_Q_std.txt: 'cell identity OK: ... hReportedCells equals the paper StatOnly set') and fig6/7 against budget_Q. The five-definition recompute is the Stage F rebuild's, on its Q (comment below).

\subsection{Edge-bin sensitivity table (from App. A, sec:chisq)}
% From app_statmethods.tex:369 @143f4fc97; table of a robustness check of the exact-split cross-check
% (\texttt{compare\_to\_paper\_fullcov.py})
\begin{center}
\begin{tabular}{@{}rrrrr@{}}
\toprule
$\pPar \ge$ (GeV/c) & $N_{\text{bins}}$ & $\chi^{2}$ & $\chi^{2}/\mathrm{ndf}$ & median ratio \\
\midrule
1.5 & 205 & 750.49 & 3.661 & 1.0064 \\
2.0 & 196 & 685.36 & 3.497 & 1.0059 \\
2.5 & 186 & 545.87 & 2.935 & 1.0059 \\
3.0 & 175 & 474.84 & 2.713 & 1.0059 \\
3.5 & 163 & 452.95 & 2.779 & 1.0073 \\
4.0 & 151 & 428.72 & 2.839 & 1.0068 \\
\bottomrule
\end{tabular}
\end{center}

\subsection{Regularization scans on an earlier covariance (from App. A, sec:rank)}
% From app_statmethods.tex:465 @143f4fc97; table of an earlier covariance and the exact-split cross-check
% NR: no committed producer of this scan (rebuild REPORT §6); kept as a labelled historical diagnostic.
% J3/J4 provisional (DECISION-20261018, packet C7: kept, labelled historical).
\begin{center}
\resizebox{\textwidth}{!}{%
\begin{tabular}{@{}lll@{}}
\toprule
Method & Setting & Result \\
\midrule
Eigenvalue truncation & $\lambda/\lambda_{\max} > \{10^{-3},10^{-4},10^{-5}\}$ &
  rank $\{19,49,82\}$; $\chi^{2}/205 = \{0.13,1.12,3.08\}$ \\
Ridge $\mathbf{C}+\varepsilon\lambda_{\max}\mathbf{I}$ & $\varepsilon=\{10^{-3},10^{-4},10^{-5}\}$ &
  $\chi^{2}/205 = \{0.35,1.38,4.30\}$ \\
Diagonal-target shrinkage & $\lambda=0.05$ &
  rank $178\to204$, cond $1.3\times10^{11}$, $\chi^{2}/205 = 3.29$ \\
\bottomrule
\end{tabular}}
\end{center}

\subsection{Multinomial to independent Poisson bootstrap limit (from App. A, app:poisson)}
% From app_statmethods.tex:126 @143f4fc97; derivation
Drawing $N$ events with replacement from an $N$-event sample gives counts $(n_{1},\dots,n_{N}) \sim \mathrm{Multinomial}(N,\mathbf{p})$ with uniform $p_{i}=1/N$ and $\sum_{i}n_{i}=N$. Each marginal $n_{i}\sim\mathrm{Binomial}(N,1/N)$ has mean $1$, variance $(N-1)/N \to 1$, and converges to $\mathrm{Poisson}(1)$ as $N$ grows (many trials, small success probability, fixed expected count). The $\sum n_{i}=N$ constraint induces $\mathcal{O}(1/N)$ correlations that vanish at large $N$. The multinomial is an independent Poisson process conditioned on $N$ counts. Removing that constraint gives independent Poisson(1) variables (Poissonization). One Poisson(1) factor is drawn per event. For $N \gtrsim 10^{6}$, our regime, the relative total-count error is $\mathcal{O}(1/\sqrt{N})$, far smaller than the per-bin Poisson variance.

\subsection{Sources of the 2D comparison numbers (from App. A, sec:assemble)}
% From app_statmethods.tex:536 @143f4fc97; an index of numbers stated in their own sections
\subsubsection*{Sources of the 2D comparison numbers}

\begin{center}
\begin{tabularx}{\textwidth}{@{}>{\RaggedRight\arraybackslash}p{0.27\textwidth}l>{\RaggedRight\arraybackslash}X@{}}
\toprule
Number & Section & Procedure \\
\midrule
$\chi^{2}/\mathrm{ndf} = \chiCombined$ & \S\ref{sec:chisq} &
  Eq.~\eqref{eq:chi2} with $\mathbf{C}^{\paper}+\mathbf{C}^{\syst}+\mathbf{C}^{\stat}+\mathbf{C}^{\ML}+\mathbf{C}^{\mathrm{PL}}$ \\
$\mathbf{C}^{\syst}$ & \S\ref{sec:universes} &
  Universe sweep ($187$ universes; $43$ bands in the band sum), MAT-conformant Eq.~\eqref{eq:matcov} $+$ norm \\
$\mathbf{C}^{\stat}$ & \S\ref{sec:bootstrap} &
  Poisson bootstrap, $300$ replicas, seed $42$, three streams, Eq.~\eqref{eq:bootcov} \\
$\mathbf{C}^{\ML}$ & \S\ref{sec:seedscan} & LightGBM training-seed settings, $5\times20$ replicas \\
$\mathbf{C}^{\mathrm{PL}}$ & \S\ref{sec:uq-combine} & Run-period block, three run periods \\
$\mathbf{C}^{\paper}$ & \S\ref{sec:minerva} & Paper's published \texttt{TotalCovariance} \\
Pull standard deviation $= 0.389$ & \S\ref{sec:pull} & % G2R: RE-2-shapes run 2 (C_new), g2r-integration/ledger.tsv g2r-audit-pull-rms-b; a source comment citing RE-2 part C for this value gives the earlier value
  Per-bin standardized residual on combined $\mathbf{C}$ \\
Combined median rel.\ $= 6.463\%$ & \S\ref{sec:combine} & % G2R: RE-2-shapes run 2 (C_new), g2r-integration/ledger.tsv g2r-audit-q-median; a source comment citing RE-2 part C for this value gives the earlier value
  $\sqrt{\diag\mathbf{C}^{\tot}}/\mathbf{X}^{\CV}$, $\mathbf{X}^{\CV}$ the LightGBM seed-42 central value \\
Fixed-truth $C_2=\covTwoDCTwo$ (nominal $0.954$) & \S\ref{sec:coverage} &
  Earlier statistical band vs.\ fixed truth, $N_{\text{toy}}=\covTwoDNtoy$: fails \\
% J3/J4 provisional (DECISION-20261018, packet W3: D2): the two coverage rows score earlier bands, not Q's statistical block (not recomputed, NR).
Same-ensemble $|r|\leq1$: $68.71\%$ & \S\ref{sec:pull-diagnostic} &
  Pull diagnostic, $N_{\text{toy}}=200$ \\
\bottomrule
\end{tabularx}
\end{center}

\subsection{Familywise error: the crossing event and the worst case (from App. A, sec:joint5d-fwer)}
% From app_statmethods_nd.tex:134 @e57f0a0a8; derivation steps
\emph{Crossing event.}  The rule rejects at threshold $c$ only if $U(k_{B},B)<c$ at the number $B\le\jtBcap$ at which sampling stopped.  Whatever rule chose that stopping point, including stops decided with results visible, irregular looks, lost draws, and the budget extensions, the rejection lies in the event $E_{c}=\bigcup_{B\le\jtBcap}\{U(k_{B},B)<c\}$ that the bound crosses below $c$ at some $B$.  Hence $P(\text{reject at }c)\le R(c)\equiv E_{p}[g_{c}(p)]$, with $g_{c}(p)=P(E_{c}\mid p)$.

\emph{Uniform $p$ is the worst case.}  Coupling the Bernoulli draws in $p$ shows that $g_{c}(p)$ does not increase with $p$.  For any super-uniform $p$ its expectation is therefore at most its value for uniform $p$, $R(c)\le\int_{0}^{1}g_{c}(u)\,du$.

\subsection{Familywise error: certified evaluation of R(c) (from App. A, sec:joint5d-fwer)}
% From app_statmethods_nd.tex:145 @e57f0a0a8; derivation steps
\emph{Certified constant.}  For each $c=\alpha/j$, $j=1,\dots,10$, let $k^{*}(B)$ be the largest $k$ with $U(k,B)<c$.  For fixed $p$, a dynamic program propagates the distribution of $k_{B}$ from $B=1$ to \jtBcap{}, moving the probability of every state with $k_{B}\le k^{*}(B)$ into an absorbing ``crossed'' state at each $B$, and pooling the counts above $\max_{B}k^{*}(B)$, which can never cross because counts never decrease. The absorbed probability is $g_{c}(p)$ exactly.

This is evaluated on a grid of 6000 points in $[0,6c]$.  Because $g_{c}$ is nonincreasing, a left Riemann sum over-estimates its integral on each grid interval, and with the tail term $(1-6c)\,g_{c}(6c)$ it is a rigorous upper bound on $\int_{0}^{1}g_{c}$, so the integration error is included in the bound.

\subsection{Same-ensemble pull diagnostic: the toy ensemble (from App. A, sec:pull-diagnostic)}
% From app_statmethods.tex:514 @e57f0a0a8; procedure
 Each toy has its own bootstrap seed $k$ and per-event Poisson(1) draws, as in production, with seeds chosen not to collide with the production replicas. Toy seeds are independent. Bins within a toy are correlated. The MC bootstrap multiplies the truth weight. It fluctuates both the unfolded result and the stored truth. The stored truth varies in all 205 reported bins, and its all-toy mean $\bar X^{\text{truth}}$ is therefore not an independent fixed truth for a split-sample coverage claim. It only centers a same-ensemble diagnostic.

\subsection{From Part 2, secobservables.tex: the paragraph beginning ``Tablereftab:observables-impl gives the implemented truth and''}
% From sec_observables.tex:6 @d55f7aad6; qualifier detail or procedure (DECISION-20261120a, second pass)
In the note, this text follows: ``$E_{\mu}$, $p_{\mu}$, and $\theta_{\mu}$ are the muon energy, momentum, and angle.''\par
$q_{0}$ is the calorimetric estimate of the energy transfer.

\subsection{From Part 2, secobservables.tex: the paragraph beginning ``The five reconstructed inputs therefore share''}
% From sec_observables.tex:41 @d55f7aad6; qualifier detail or procedure (DECISION-20261120a, second pass)
In the note, this text follows: ``The five reconstructed inputs therefore share information and do not, by themselves, provide five independent constraints: the result is five simultaneous observables, not five independent measurements.''\par
The five classifier inputs, the lower-dimensional support they populate, and the five reported observables are distinct. Only reconstructed $W$ is fixed by the other inputs; reconstructed $\Eavail$ comes from a separate calorimetric estimator and is not a function of them. The nearly degenerate correlation directions of the joint-test metric (\S\ref{sec:joint5d-design}) are consistent with this redundancy, but that origin has not been shown, and the information added by $W$ has not been measured.
% carried with the text from sec_observables.tex:42 @d55f7aad6
% absence-ok: "has not been shown" and "has not been measured" are the base's own limitations of the five inputs, reworded in r22 Step 2

\subsection{From Part 2, secobservables.tex: the paragraph beginning ``The analysis excludes events with negative''}
% From sec_observables.tex:44 @d55f7aad6; qualifier detail or procedure (DECISION-20261120a, second pass)
In the note, this text follows: ``The generators measure $\pT$ and $\pPar$ about their own neutrino axis; their $\pT$ differs from the analysis truth $\pT$ by \SI{\pm4}{MeV} for 68\% of events.''\par
Over the 109 joint-test cells, this convention changes the analysis simulation's cell contents by factors between 0.9965 and 1.0035 (median 1.000008), at the level of its Monte Carlo precision, and shifts GENIE's central prediction by between $-0.47$ and $+0.53$ of each cell's null standard deviation.
% carried with the text from sec_observables.tex:45 @d55f7aad6
% Search: information added by W; W redundancy; near-degenerate directions; origin; joint-test metric; pT axis; neutrino axis; beam axis; MP-5; R_c (none for W's information or the origin: M2-ARCHIVE-CHECKS.md:7, r12/ledger.md:103-104 restate it; the pT-axis effect is printed from MP-5, state/measurement-paper/phase1/mp5/REPORT.md:66, :68, :100, job 60017887, with its transfer to the other generators UNRELIABLE, :87-90; Joseph, DECISION-20261027-joseph-absence-audit-items-3-4.md item 1) @ 7a6ae651

\subsection{From Part 2, secexperiment.tex: the paragraph beginning ``MINERvAciteMINERvA:2013zvz was a fine-grained scintillator-tracker neutrino''}
% From sec_experiment.tex:5 @d55f7aad6; qualifier detail or procedure (DECISION-20261120a, second pass)
In the note, this text follows: ``MINERvA~\cite{MINERvA:2013zvz} was a fine-grained scintillator-tracker neutrino experiment on Fermilab's NuMI beamline.''\par
Its central tracker comprises hexagonal planes of triangular polystyrene scintillator strips in three stereo orientations, surrounded by electromagnetic and hadronic calorimetry. The magnetized MINOS near detector, \SI{2}{m} downstream, serves as the muon spectrometer, measuring momentum and charge for MINOS-matched tracks.

\subsection{From Part 2, secexperiment.tex: the paragraph beginning ``The reference simulation uses GENIE 2.12.6citeAndreopoulos:2009rq''}
% From sec_experiment.tex:10 @d55f7aad6; qualifier detail or procedure (DECISION-20261120a, second pass)
In the note, this text follows: ``The reference simulation uses GENIE 2.12.6~\cite{Andreopoulos:2009rq} and MINERvA Tune v1, applied as per-event weights by the reweighting tools of the MINERvA Analysis Toolkit (MAT).''\par
These are the flux-and-central-value (flux-and-CV) weight, the GENIE model-parameter (``dial'') weights, the low-recoil 2p2h enhancement fit, and the non-resonant-pion reduction.

\subsection{From Part 2, secexperiment.tex: the paragraph beginning ``Genuine backgrounds are selected events whose''}
% From sec_experiment.tex:26 @d55f7aad6; qualifier detail or procedure (DECISION-20261120a, second pass)
In the note, this text follows: ``Out-of-phase-space signal, reconstructed inside the phase space with its true muon outside it, is routed into the background template with them.''\par
Both are POT-scaled and subtracted from the reconstructed-data target before unfolding (\S\ref{sec:method}). Their sum, $N_{\mathrm{bkg}}(b)$ in each reconstructed bin, enters the purity weight $w_{\mathrm{pur}}(b)$ of \S\ref{sec:method-algo}.

\subsection{From Part 2, secexperiment.tex: the paragraph beginning ``The background template used for this''}
% From sec_experiment.tex:28 @d55f7aad6; qualifier detail or procedure (DECISION-20261120a, second pass)
In the note, this text follows: ``The background template used for this result is \SI{2.91}{\percent} of selected data events in the two-dimensional grid over the full sample (\SIrange{2.72}{3.23}{\percent} per playlist, \SI{2.74}{\percent} on playlist 1A).''\par
On playlist 1A only, about \SI{10}{\percent} (weighted) of the template is out-of-phase-space signal. These are $\nu_{\mu}$ charged-current events with a true vertex inside the fiducial volume. The rest is genuine background, mostly $\nu_{\mu}$ charged-current events with a true vertex outside it.

\subsection{From Part 2, secexperiment.tex: the paragraph beginning ``Selected data exceed the POT-scaled MINERvA''}
% From sec_experiment.tex:39 @d55f7aad6; qualifier detail or procedure (DECISION-20261120a, second pass)
Selected data exceed the POT-scaled MINERvA Tune~v1 prediction (signal $+$ background) by $\sim\SI{12}{\percent}$ in overall normalization.

\subsection{From Part 2, secexperiment.tex: the paragraph beginning ``Selected data exceed the POT-scaled MINERvA''}
% From sec_experiment.tex:39 @d55f7aad6; qualifier detail or procedure (DECISION-20261120a, second pass)
In the note, this text follows: ``The data/MC ratio rises with $\pT$, from $0.88$ in the lowest bin to $\sim1.2$ above \SI{1}{GeV/c}.''\par
Excess and trend together reflect the truth-level normalization deficit of \S\ref{sec:results-model} and \S\ref{sec:3d-models}. Beyond the $\pT$ trend noted above, the ratio panels also show a mild enhancement in the lowest $\Eavail$ bins, reflecting the low-recoil excess examined at truth level in \S\ref{sec:3d-2p2h}.

\subsection{From Part 2, secmethod.tex: the paragraph beginning ``item textbfStep 1 (detector space). A''}
% From sec_method.tex:9 @d55f7aad6; qualifier detail or procedure (DECISION-20261120a, second pass)
In the note, this text follows: ``The weight is \emph{pulled back} directly to its truth record.''\par
Truth events without a reconstructed partner (``misses'') receive step-1 weights from a regressor trained on the reconstructed subsample, $w(\mathrm{truth\ features})$. This native OmniFold acceptance treatment requires no separate efficiency correction.

\subsection{From Part 2, secmethod.tex: the paragraph beginning ``Unbinned OmniFold adds an observable by''}
% From sec_method.tex:14 @d55f7aad6; qualifier detail or procedure (DECISION-20261120a, second pass)
Unbinned OmniFold adds an observable by adding an input to the step-1 and step-2 classifiers, without an explicit migration matrix or per-bin response tensor. It avoids the combinatorial binning penalty of matrix unfolding but still has a dimensional cost: classifier training, acceptance and support coverage, sample sparsity and uncertainty quantification become more demanding as observables are added. The primary methodological motivation is the absence of a practical analog of D'Agostini iterative Bayesian unfolding (IBU).

\subsection{From Part 2, secmethod.tex: the paragraph beginning ``paragraphBackground subtraction. In binned purity subtraction,''}
% From sec_method.tex:21 @d55f7aad6; qualifier detail or procedure (DECISION-20261120a, second pass)
No negative-weight events are injected.

\subsection{From Part 2, secmethod.tex: the paragraph beginning ``For other experiments' subtraction methods, see''}
% From sec_method.tex:32 @d55f7aad6; qualifier detail or procedure (DECISION-20261120a, second pass)
For other experiments' subtraction methods, see \S\ref{sec:intro}.

\subsection{From Part 2, secmethod.tex: the paragraph beginning ``For other experiments' subtraction methods, see''}
% From sec_method.tex:32 @d55f7aad6; qualifier detail or procedure (DECISION-20261120a, second pass)
In the note, this text follows: ``A fully unbinned alternative adds negative-weight background simulation on the measured side.''\par
Its bin-integrated target is $(D_{b}-B_{b})/S_{b}$ (data, background and signal).

\subsection{From Part 2, secmethod.tex: the paragraph beginning ``For other experiments' subtraction methods, see''}
% From sec_method.tex:32 @d55f7aad6; qualifier detail or procedure (DECISION-20261120a, second pass)
In the note, this text follows: ``The alternative reproduces the binned correction to \SI{\nwPctTot}{\percent} on the total two-dimensional (2D) cross section and its systematic and statistical covariances to within \SI{2}{\percent} (subtraction comparison, App.~\ref{app:negweight}).''\par
It removes residual binning dependence and suits the extended-fiducial analysis with larger backgrounds. The quoted 2D result and covariance, 3D and 4D central values, sixteen-cell result and covariance, and studies and tests of the \estLowCap{} and \estMultiD{} use the refined variant. All other quoted cross sections and uncertainties, including the \resFiveAxis{} of the \estPurity{} and the extended-fiducial value, use binned subtraction. The matrix, Table~\ref{tab:estimators-matrix} (configurations, \S\ref{sec:method-config}), links results to configurations.

\subsection{From Part 2, secmethod.tex: the paragraph beginning ``The quoted 2D cross section and''}
% From sec_method.tex:48 @d55f7aad6; qualifier detail or procedure (DECISION-20261120a, second pass)
In the note, this text follows: ``The quoted 2D cross section and covariance use the \estTwoD{}.''\par
The 3D and 4D central values, the \resSixteen{} and the generator tests use the \estMultiD{}. Both subtract background with negative weights and a 400-tree refinement classifier. They differ in training histogram construction in the LightGBM learner and in deterministic settings. An earlier \estPurity{} is retained for the \resFiveAxis{} and for specified figures and closures. Three other configurations serve only as diagnostic controls in the studies listed in Table~\ref{tab:estimators-matrix}, the matrix of estimator uses and implementations.

\subsection{From Part 2, secmethod.tex: the paragraph beginning ``Gradient-boosted decision trees (GBDTs) replace OmniFold's''}
% From sec_method.tex:50 @d55f7aad6; qualifier detail or procedure (DECISION-20261120a, second pass)
Gradient-boosted decision trees (GBDTs) replace OmniFold's original neural networks when unfolding scalar observables. With at most five high-level features per event, tree ensembles typically match or outperform neural networks~\cite{Grinsztajn:2022tabular} and train deterministically and far faster on a central processing unit (CPU). All unfolding results here use GBDTs (point-cloud representations, App.~\ref{sec:pet}). Calibration tests and multilayer perceptron (MLP) cross-checks test the implementation (classifier checks, \S\ref{sec:3d-classifier}).

\subsection{From Part 2, secmethod.tex: the paragraph beginning ``The unfolding code supports interchangeable learners''}
% From sec_method.tex:52 @d55f7aad6; qualifier detail or procedure (DECISION-20261120a, second pass)
The unfolding code supports interchangeable learners (component terminology and code paths in App.~\ref{app:codebase}). The quoted 2D central value and all its uncertainty ensembles use LightGBM with refined negative-weight subtraction (background method, App.~\ref{app:negweight}). The estimator matrix (Table~\ref{tab:estimators-matrix}) lists this as the \estTwoD{}. The ensembles cover the systematic-universe sweep, Poisson bootstrap, training-seed variations and run-period unfoldings.

\subsection{From Part 2, secmethod.tex: the paragraph beginning ``The unfolding code supports interchangeable learners''}
% From sec_method.tex:52 @d55f7aad6; qualifier detail or procedure (DECISION-20261120a, second pass)
In the note, this text follows: ``The step-1 classifier, the step-2 classifier and the regressor for unreconstructed truth events share one setting: $100$ trees, $8$ leaves and learning rate $0.1$.''\par
The universe sweep supplies the central value at seed~42. Thus the systematic covariance shares its estimator and seed (systematic universes, \S\ref{sec:uq-universes}). The quoted statistical covariance also uses seed~42, with three independent bootstrap streams (data, signal and background simulation, \S\ref{sec:uq-boot}). The training-seed covariance uses five seed settings that include seed~42 (\S\ref{sec:uq-ml}).

\subsection{From Part 2, secmethod.tex: the paragraph beginning ``A central unfolding at seed45 differs''}
% From sec_method.tex:54 @d55f7aad6; qualifier detail or procedure (DECISION-20261120a, second pass)
In the note, this text follows: ``A central unfolding at seed~45 differs from the seed-42 one by at most \SI{1.4}{\percent} in a cell.''\par
The quoted central value has a recorded configuration and repeats numerically. It agrees with the refined 2D data unfolding (background study, App.~\ref{sec:negweight-footing}).
% carried with the text from sec_method.tex:55 @d55f7aad6
% Source: docs/publication/DECISION-20261011-joseph-2d-central-lightgbm-seed42.md; receipts

\subsection{From Part 2, secmethod.tex: the paragraph beginning ``A central unfolding at seed45 differs''}
% From sec_method.tex:54 @d55f7aad6; qualifier detail or procedure (DECISION-20261120a, second pass)
In the note, this text follows: ``The purity-subtracted seed-42 central value has no recorded run configuration, although an independent rerun and a regeneration whose run configuration is recorded both reproduce it (Ref.~\cite{NoteRecords2026}, \recentry{2d-central-value}).''\par
Against the ten purity-subtracted LightGBM central values at seeds 1--10 it is mildly over-dispersed (root-mean-square standardized difference 1.28, maximum 4.25). The choice of seed~42 is recorded in App.~\ref{app:devhistory}.
% carried with the text from sec_method.tex:55 @d55f7aad6
% Source: docs/publication/DECISION-20261011-joseph-2d-central-lightgbm-seed42.md; receipts

\subsection{From Part 2, secmethod.tex: the paragraph beginning ``paragraphThe exact-split cross-check. The exact-split configuration''}
% From sec_method.tex:85 @d55f7aad6; qualifier detail or procedure (DECISION-20261120a, second pass)
Its seed and run log appear in the development history (App.~\ref{app:devhistory}). The result depends on the estimator seed. A reseeded run (\texttt{random\_state} 1/2/3) moves every bin (median \SI{0.18}{\percent}, maximum \SI{1.6}{\percent}) and gives $\chi^{2}/\mathrm{ndf}=\chiPaperExactSeeded$ (chi-square per degree of freedom). The spread over seeds is unmeasured.
% carried with the text from sec_method.tex:86 @d55f7aad6
% Search: exact-split seed scan; spread over seeds; matched covariance; exact-split bootstrap; P09b @ ddf16f10
% Search: exact-backend seed scan; exact-split seed scan; P09b (pairings.tsv: "no exact-backend seed scan exists") @ 1a015eef

\subsection{From Part 2, secmethod.tex: the paragraph beginning ``paragraphThe exact-split cross-check. The exact-split configuration''}
% From sec_method.tex:85 @d55f7aad6; qualifier detail or procedure (DECISION-20261120a, second pass)
In the note, this text follows: ``The cross-check has no matched covariance.''\par
No exact-split bootstrap, universe sweep or seed scan exists.
% carried with the text from sec_method.tex:86 @d55f7aad6
% Search: exact-split seed scan; spread over seeds; matched covariance; exact-split bootstrap; P09b @ ddf16f10
% Search: exact-backend seed scan; exact-split seed scan; P09b (pairings.tsv: "no exact-backend seed scan exists") @ 1a015eef


% carried with the text from sec_method.tex:86 @d55f7aad6
% Search: exact-split seed scan; spread over seeds; matched covariance; exact-split bootstrap; P09b @ ddf16f10
% Search: exact-backend seed scan; exact-split seed scan; P09b (pairings.tsv: "no exact-backend seed scan exists") @ 1a015eef

\subsection{From Part 2, secmethod.tex: the paragraph beginning ``paragraphThe estimator configurations. The six estimator''}
% From sec_method.tex:102 @d55f7aad6; qualifier detail or procedure (DECISION-20261120a, second pass)
Each result names its estimator.

\subsection{From Part 2, secmethod.tex: the paragraph beginning ``paragraphThe estimator configurations. The six estimator''}
% From sec_method.tex:102 @d55f7aad6; qualifier detail or procedure (DECISION-20261120a, second pass)
``Deterministic'' means fixed seeds and reproducible training.

\subsection{From Part 2, secmethod.tex: the paragraph beginning ``paragraphThe estimator configurations. The six estimator''}
% From sec_method.tex:102 @d55f7aad6; qualifier detail or procedure (DECISION-20261120a, second pass)
In the note, this text follows: ``Results from different configurations are not combined.''\par
The \estTwoD{} adds refined negative-weight subtraction to the nominal LightGBM configuration above. It uses the \estMultiD{} refinement classifier at the estimator seed plus three. Its quoted seed-42 central value repeats numerically (above). The \estPurity{} is the same learner with purity-based background subtraction (algorithm, \S\ref{sec:method-algo}). The \estMultiD{}, the estimator of the generator tests, is a deterministic LightGBM variant with the unbinned negative-weight background treatment (background method, App.~\ref{app:negweight}). Its refinement classifier has 400 trees and 31 leaves.

\subsection{From Part 2, secmethod.tex: the paragraph beginning ``paragraphThe estimator configurations. The six estimator''}
% From sec_method.tex:102 @d55f7aad6; qualifier detail or procedure (DECISION-20261120a, second pass)
In the note, this text follows: ``The negative-weight treatment reduces the purity method's nominal-truth bias in the highest-$W$ (hadronic invariant mass) column by about an order of magnitude (from 3.5--\SI{4.1}{\percent} with purity subtraction to 0.11--\SI{0.24}{\percent} with negative weights).''\par
This reaches the no-background closure level. The larger refinement classifier removes most of the remaining background-related bias (background and departure tests, \S\S\ref{sec:val-negweight}--\ref{sec:val-departure}). The \estMultiD{} was selected before the joint-test design was fixed (\S\ref{sec:joint5d}). The \estExact{} is described above.

\subsection{From Part 2, secmethod.tex: the paragraph beginning ``paragraphHyperparameter choice. Neither OmniFold classifier nor''}
% From sec_method.tex:144 @d55f7aad6; qualifier detail or procedure (DECISION-20261120a, second pass)
The \estMultiD{} keeps five iterations under the precision study's rule (\S\ref{sec:gbdt-model-dependence}). Its refinement classifier, 400 trees and 31 leaves, was selected before the joint-test design was fixed (\S\ref{sec:joint5d}).

\subsection{From Part 2, secmethod.tex: the paragraph beginning ``The fixed settings are the scikit-learn''}
% From sec_method.tex:146 @d55f7aad6; qualifier detail or procedure (DECISION-20261120a, second pass)
The fixed settings are the scikit-learn \texttt{GradientBoosting} defaults: 100 trees, learning rate $0.1$, and depth-3 capacity. The LightGBM learner of the central value and its ensembles uses matched capacity ($\texttt{num\_leaves}=8=2^{3}$), so both have the same deliberately modest capacity, an OmniFold regularizer (algorithm, \S\ref{sec:method-algo}).

\subsection{From Part 2, secmethod.tex: the paragraph beginning ``The estPurity and the estTwoD leave''}
% From sec_method.tex:148 @d55f7aad6; qualifier detail or procedure (DECISION-20261120a, second pass)
The \estPurity{} and the \estTwoD{} leave histogram construction at its default. Bin edges come from the library's default row sample. The training seed affects the result through that sample. The configurations this section calls deterministic set those edges from every training row and enable the library's deterministic training.

\subsection{From Part 2, secmethod.tex: the paragraph beginning ``The Keras cross-check (Srefsec:3d-classifier) uses the''}
% From sec_method.tex:150 @d55f7aad6; qualifier detail or procedure (DECISION-20261120a, second pass)
The Keras cross-check (\S\ref{sec:3d-classifier}) uses the same multilayer perceptron for both classifiers and missed-event regression. It has three hidden layers of 64, 128 and 64 units and Gaussian error linear unit (GELU) activation. Training uses Adam for at most 50 epochs with early stopping, per-fit input standardization and a fixed shuffle. Settings stay fixed. Agreement does not address sensitivity to single-precision (float32) ties (departure tests, \S\ref{sec:val-departure}).

\subsection{From Part 2, secmethod.tex: the paragraph beginning ``The following diagnostics report performance. They''}
% From sec_method.tex:152 @d55f7aad6; qualifier detail or procedure (DECISION-20261120a, second pass)
In the note, this text follows: ``The following diagnostics report performance. They do not validate the quoted estimator or its intervals. They include closure tests (\S\ref{sec:validation}) and classifier calibration (\S\ref{sec:3d-classifier}).''\par
The exact-split cross-check measures insensitivity to doubling the iterations (below). Method-dependence studies report backend agreement (\S\ref{sec:results-tension}). The training-seed variation band measures residual machine-learning (ML) stochasticity (\S\ref{sec:systematics}).

\subsection{From Part 2, secmethod.tex: the paragraph beginning ``Production uses n=5 iterations. Between five''}
% From sec_method.tex:154 @d55f7aad6; qualifier detail or procedure (DECISION-20261120a, second pass)
Production uses $n=5$ iterations.

\subsection{From Part 2, secmethod.tex: the paragraph beginning ``Production uses n=5 iterations. Between five''}
% From sec_method.tex:154 @d55f7aad6; qualifier detail or procedure (DECISION-20261120a, second pass)
In the note, this text follows: ``Between five and ten iterations, the exact-split total cross section agrees to within \SI{0.026}{\percent} (validation, \S\ref{sec:validation}). That measures movement of this diagnostic for that cross-check. It does not establish convergence of every functional.''\par
Separately, the unfolded result reproduces the integrated cross section of an independently published binned measurement of the same data (results, \S\ref{sec:results}). That comparison is not an iteration test.

\subsection{From Part 2, secmethod.tex: the paragraph beginning ``Every production, universe-sweep and closure run''}
% From sec_method.tex:164 @d55f7aad6; qualifier detail or procedure (DECISION-20261120a, second pass)
In the note, this text follows: ``Every production, universe-sweep and closure run follows the same pipeline conventions (records document, Ref.~\cite{NoteRecords2026}).''\par
Only events passing reconstructed-level selection reach the step-1 classifier.

\subsection{From Part 2, secmethod.tex: the paragraph beginning ``Every production, universe-sweep and closure run''}
% From sec_method.tex:164 @d55f7aad6; qualifier detail or procedure (DECISION-20261120a, second pass)
In the note, this text follows: ``Both sides of step-1 training exclude events with truth outside the measurement phase space.''\par
Closure pseudo-data follow the production data path.

\subsection{From Part 2, secmethod.tex: the paragraph beginning ``In n dimensions, the differential cross''}
% From sec_method.tex:176 @d55f7aad6; qualifier detail or procedure (DECISION-20261120a, second pass)
In the note, this text follows: ``Thus the defined truth sample enters and $c\equiv1$ per bin.''\par
Negative truth $W^{2}$ is excluded by the five-observable sample definition (\S\ref{sec:eavailw-definitions}). This differs from losing a reconstructed event from the input. The factor remains a monitored self-check. Dropping unmatched truth events would give $c<1$.

\subsection{From Part 2, secmethod.tex: the paragraph beginning ``In n dimensions, the differential cross''}
% From sec_method.tex:176 @d55f7aad6; qualifier detail or procedure (DECISION-20261120a, second pass)
In the note, this text follows: ``$\Phi$ is the flux integral by transverse muon momentum $\pT$. It is tabulated per square meter and broadcast over other axes.''\par
The formula uses that value.

\subsection{From Part 2, secmethod.tex: the paragraph beginning ``In n dimensions, the differential cross''}
% From sec_method.tex:176 @d55f7aad6; qualifier detail or procedure (DECISION-20261120a, second pass)
In the note, this text follows: ``The factor $10^{4}$ converts flux tabulated per square meter to the reported $\si{cm^2}$ ($\SI{1}{m^2}=10^{4}\,\si{cm^2}$), giving units of $\si{cm^2}$ per axis unit per nucleon.''\par
The implementation (code, App.~\ref{app:codebase}) is self-tested to reproduce stored lower-dimensional reference extractions to $<10^{-12}$. Since $\sigma\propto1/\Phi$, $\Phi$ varies per flux universe when building the systematic covariance. This variation dominates the Flux band (systematics, \S\ref{sec:systematics}).

\subsection{From Part 2, secmethod.tex: the paragraph beginning ``For higher-dimensional extractions the same integral''}
% From sec_method.tex:179 @d55f7aad6; qualifier detail or procedure (DECISION-20261120a, second pass)
In the note, this text follows: ``For higher-dimensional extractions the same integral is $\SI{6.243e-8}{cm^{-2}}$ per proton on target, from 0 to \SI{100}{GeV} after the neutrino--electron-scattering constraint.''\par
$N$ is the formula's fiducial nucleon count. The $\si{cm^{-2}}$ value is the converted integral. It is not substituted for $\Phi$.

\subsection{From Part 2, secmethod.tex: the paragraph beginning ``The 2D grid has 14 pT''}
% From sec_method.tex:184 @d55f7aad6; qualifier detail or procedure (DECISION-20261120a, second pass)
In the note, this text follows: ``The 19 unreported diagonal cells have $\pT^{\mathrm{hi}}/\pPar^{\mathrm{lo}}>\tan20^{\circ}$.''\par
The extension axes use $\Eavail$: $[0,0.1,0.2,0.4,0.8,1.5,3.0,100]\,\si{GeV}$; $q_{3}$ (three-momentum transfer): $[0,0.2,0.4,0.6,0.8,1.2,2.0,100]\,\si{GeV}$; $W$: $[0,1.1,1.4,1.8,2.2,3.0,100]\,\si{GeV}$.

\subsection{From Part 2, secmethod.tex: the paragraph beginning ``The 2D grid has 14 pT''}
% From sec_method.tex:184 @d55f7aad6; qualifier detail or procedure (DECISION-20261120a, second pass)
In the note, this text follows: ``Wide final bins capture full charged-current (CC) inclusive tails in marginals.''\par
Truncating the top edge would break the marginalization checks in \S\ref{sec:3d-anchor}. Post-unfolding binning lets us revisit these choices without rerunning unfolding.

\subsection{From Part 2, secmethod.tex: the paragraph beginning ``Truthtoreco migration on the analysis binning''}
% From sec_method.tex:187 @d55f7aad6; qualifier detail or procedure (DECISION-20261120a, second pass)
Truth$\to$reco migration on the analysis binning (Fig.~\ref{fig:migration}) shows the smearing OmniFold must undo. OmniFold uses no response matrix. Median reconstruction stability (same-bin fraction, Fig.~\ref{fig:migration}) is $\sim0.6$ on each axis. Such migrations would badly bias a bin-by-bin correction. Unfolding is required.

\subsection{From Part 2, secmethod.tex: the paragraph beginning ``Profile OmniFoldciteZhu:2025profile extends classifier-based expectation maximization''}
% From sec_method.tex:199 @d55f7aad6; qualifier detail or procedure (DECISION-20261120a, second pass)
Profile OmniFold~\cite{Zhu:2025profile} extends classifier-based expectation maximization (EM) to forward-model nuisance parameters. Profiling nuisances with observed data would change our inference model, which propagates an external nuisance ensemble through a fixed procedure. Neither a profiled estimator nor its uncertainty enters here.

\subsection{From Part 2, secmethod.tex: the paragraph beginning ``SPINUP learns a forward mappingciteButter:2025spinup. AUSSIE''}
% From sec_method.tex:201 @d55f7aad6; qualifier detail or procedure (DECISION-20261120a, second pass)
SPINUP learns a forward mapping~\cite{Butter:2025spinup}. AUSSIE replaces OmniFold's second step with a different loss and avoids the outer iteration~\cite{Ore:2026aussie}. Both offer possible future cross-checks with different algorithmic assumptions.

\subsection{From Part 2, secmethod.tex: the paragraph beginning ``SPINUP learns a forward mappingciteButter:2025spinup. AUSSIE''}
% From sec_method.tex:201 @d55f7aad6; qualifier detail or procedure (DECISION-20261120a, second pass)
In the note, this text follows: ``Reduced simulation-prior dependence does not establish detector-model independence, adequate phase-space support, or coverage for this result.''\par
Comparisons would require matched signal, selection, normalization, and uncertainty definitions.

\subsection{From App. A, appstatmethods.tex: the paragraph beginning ``This appendix gives the statistical methods''}
% From app_statmethods.tex:5 @d55f7aad6; qualifier detail or procedure (DECISION-20261120a, second pass)
This appendix gives the statistical methods of this note's results, centered on the 2D comparison: covariance inversion; the training-seed, bootstrap and universe estimators, with the Poisson limit and the $\pm1\sigma$ pair covariance; their MAT-conformant construction and combination; the $\chi^{2}$, pull and shape-only conventions; the rank and conditioning of the covariances; the anatomy of the tension in the 2D comparison; the closure procedures; and a summary of the 2D comparison numbers. It covers analysis-specific conventions, not general statistics. The $\chi^{2}$ convention used throughout is Eq.~\eqref{eq:chi2} in \S\ref{sec:chisq}.

\subsection{From App. A, appstatmethods.tex: the paragraph beginning ``A sample covariance from Nle n''}
% From app_statmethods.tex:20 @d55f7aad6; qualifier detail or procedure (DECISION-20261120a, second pass)
A sample covariance from $N\le n$ replicas in $n$ bins has rank at most $N-1$.

\subsection{From App. A, appstatmethods.tex: the paragraph beginning ``A sample covariance from Nle n''}
% From app_statmethods.tex:20 @d55f7aad6; qualifier detail or procedure (DECISION-20261120a, second pass)
In the note, this text follows: ``The 300-replica bootstrap ($N=300 > 205$) has more replicas than bins, which permits full rank and does not establish it.''\par
The rank of the systematic block is limited at any universe count (\S\ref{sec:rank}).

\subsection{From App. A, appstatmethods.tex: the paragraph beginning ``A coherent bootstrap carries each data''}
% From app_statmethods.tex:103 @d55f7aad6; qualifier detail or procedure (DECISION-20261120a, second pass)
In the note, this text follows: ``A coherent bootstrap carries each data factor through step-1 classifier training and final cross-section construction.''\par
Each MC factor is shared by the reco-tree and truth-tree copies, preserving migration consistency. The three sub-RNGs (data, signal MC, background MC; Table~\ref{tab:variance-sources}) make the sources uncorrelated by construction, and pinning the GBDT seed isolates the bootstrap from the ML noise of \S\ref{sec:seedscan}. A Poisson(1) factor can be zero. Dropping those events from each replica mimics a resample with replacement. Poisson(1) is the large-$N$ limit of the multinomial counts (\S\ref{app:poisson}).

\subsection{From App. A, appstatmethods.tex: the paragraph beginning ``For these replicas (205 bins, MEFHC''}
% From app_statmethods.tex:117 @d55f7aad6; qualifier detail or procedure (DECISION-20261120a, second pass)
In the note, this text follows: ``For these replicas (205 bins, ME~FHC LightGBM, seed~42) the relative standard deviation of the total cross section is $\boxed{0.061\%}$, the per-bin median relative spread $\boxed{0.676\%}$, and the per-bin $p_{84}$ $1.602\%$.''\par
With $300$ replicas the per-bin $\sigma$ has a relative standard error of about $4.1\%$ by the Gaussian formula $1/\sqrt{2(N-1)}$. $N=300$ is chosen for full rank.

\subsection{From App. A, appstatmethods.tex: the paragraph beginning ``The universe sweep used by the''}
% From app_statmethods.tex:130 @d55f7aad6; qualifier detail or procedure (DECISION-20261120a, second pass)
In the note, this text follows: ``The universe sweep used by the quoted budget has 187 universes across 44 bands, dominated by the $100$-universe PPFX flux band (hadron-production sampling).''\par
The quoted band sum drops one band (Rvp1pi) and holds $43$ bands plus the normalization.

\subsection{From App. A, appstatmethods.tex: the paragraph beginning ``The universe sweep used by the''}
% From app_statmethods.tex:130 @d55f7aad6; qualifier detail or procedure (DECISION-20261120a, second pass)
In the note, this text follows: ``Almost all other bands are $\pm1\sigma$ pairs ($N_{u}=2$).''\par
They include the MINOS and lateral MINERvA muon energy scales, the MINOS muon-match efficiency, the GENIE parameters \texttt{MaRES}, \texttt{MvRES} and \texttt{MaCCQE}, and the lateral GEANT FSI responses for neutrons, pions and protons. Where selection migration matters, the five kinematic bands (the two beam-angle components, the muon resolution and the two muon energy scales) require re-selected event loops, run with one universe active at a time (weight-only bands: \S\ref{sec:uq-higherdim}). The quoted universe block takes these five bands from re-unfolds of re-selected universes.

\subsection{From App. A, appstatmethods.tex: the paragraph beginning ``MINERvA's MAT framework (textttPlotUtils/allowbreak MnvVertErrorBand::allowbreak CalcCovMx,''}
% From app_statmethods.tex:146 @d55f7aad6; qualifier detail or procedure (DECISION-20261120a, second pass)
In the note, this text follows: ``uniformly for all $N_{u}\ge 2$, with no special $\pm1\sigma$ pair formula.''\par
At $N_{u}=2$ it reduces to rank-1 $\tfrac14(\mathbf{X}_{+}-\mathbf{X}_{-})(\mathbf{X}_{+}-\mathbf{X}_{-})\T$. Mean-centering removes a common offset from the central value. The half-difference remains and need not be small. For a one-sided or asymmetric pair, CV-centering adds that offset as a second direction (\S\ref{app:pair}).

\subsection{From App. A, appstatmethods.tex: the paragraph beginning ``Equationeqrefeq:systblocksum omits a larger correction. The''}
% From app_statmethods.tex:159 @d55f7aad6; qualifier detail or procedure (DECISION-20261120a, second pass)
In the note, this text follows: ``A unified throw measures it through $N$ re-unfolds, simultaneously perturbing every reweightable band and the flux multisim per event ($w_{\CV}\prod_B\rho_B^{g_B}$, $g_B\sim\mathcal N(0,1)$), then re-extracting the cross section.''\par
The corrected construction follows. Its use in the five-axis covariance of the \resFiveAxis{} is in \S\ref{sec:eavailw-adopted} (Eq.~\eqref{eq:zinflation}).
In the note, the corrected construction that follows is the unified throw, described in the next paragraph of App.~\ref{app:statmethods}.\par

\subsection{From App. A, appstatmethods.tex: the paragraph beginning ``Grouped by category as in Fig.reffig:uqbands,''}
% From app_statmethods.tex:199 @d55f7aad6; qualifier detail or procedure (DECISION-20261120a, second pass)
In the note, this text follows: ``Grouped by category as in Fig.~\ref{fig:uqbands}, which follows Figs.~6--7 of Ref.~\cite{MINERvA:2021owq}, the medians are: Muon reconstruction $3.437\%$, Models $1.428\%$, Flux $3.723\%$, Hadronic response $0.752\%$, Statistical $0.676\%$, Normalization $1.400\%$, Training seed $0.048\%$, Run period $2.087\%$.''\par
(MINOS$+$MINERvA) and the MINOS-efficiency band dominate; the first is the band that the Ruterbories paper identifies as flux-correlated (\S\ref{sec:rank}).

\subsection{From App. A, appstatmethods.tex: the paragraph beginning ``For two universes u=0,1 at pm''}
% From app_statmethods.tex:227 @d55f7aad6; qualifier detail or procedure (DECISION-20261120a, second pass)
The estimators agree only for CV-symmetric pairs ($\mathbf{m}=0$), both reducing to $\mathbf{a}\mathbf{a}\T$. For asymmetric or one-sided pairs, CV-centering adds the rank-1 $\mathbf{m}\mathbf{m}\T$ common-offset direction. Every per-band block $C_{b}$ here is mean-centered to match MAT (\S\ref{sec:minerva}); this matters for BeamAngle, MuonResolution and the one-sided NC band.

\subsection{From App. A, appstatmethods.tex: the paragraph beginning ``MAT's C++ framework defines the covariance''}
% From app_statmethods.tex:233 @d55f7aad6; qualifier detail or procedure (DECISION-20261120a, second pass)
MAT's C++ framework defines the covariance construction that the published paper leaves unspecified. It is the mean-centered $1/N$ per-band sample covariance and element-wise band sum used here.

\subsection{From App. A, appstatmethods.tex: the paragraph beginning ``textttMnvVertErrorBand::CalcCovMx implements Eq.eqrefeq:matcov, referenced to the''}
% From app_statmethods.tex:236 @d55f7aad6; qualifier detail or procedure (DECISION-20261120a, second pass)
\texttt{MnvVertErrorBand::CalcCovMx} implements Eq.~\eqref{eq:matcov}, referenced to the universe mean $\langle x\rangle_{u}$ rather than $X_{\CV}$, normalized by biased $1/N$ rather than Bessel's $1/(N-1)$ (a $1\%$ offset for PPFX, $N=100$), with no special case for $N=2$. The conventional CV-centered pair formula (``MINERvA-101 $\pm1\sigma$ pair''), $\tfrac12(\boldsymbol{\Delta}_{+}\boldsymbol{\Delta}_{+}\T + \boldsymbol{\Delta}_{-}\boldsymbol{\Delta}_{-}\T)$, agrees with it only for CV-symmetric shifts. The quoted universe block, filled into a \texttt{PlotUtils::MnvH1D} with MAT's universe weights on Flux, agrees with \texttt{MnvH1D::GetTotalErrorMatrix} to $1.3\times10^{-16}$ of its largest element. The quoted Flux band carries MAT's $\nu$--$e$ and inverse-muon-decay constraint weights.

\subsection{From App. A, appstatmethods.tex: the paragraph beginning ``textttMnvH1D::GetTotalErrorMatrix sums bands and the statistical''}
% From app_statmethods.tex:240 @d55f7aad6; qualifier detail or procedure (DECISION-20261120a, second pass)
In the note, this text follows: ``\texttt{MnvH1D::GetTotalErrorMatrix} sums bands and the statistical matrix element-wise, assuming band independence.''\par
Cross-band correlations exist only if encoded in a combined band, such as the Bashyal joint fit handled outside this loop. MAT applies no ridge, shrinkage, SVD truncation or rank-deficiency handling to the published total covariance. The statistical block makes $V$ full-rank. Its systematic-only rank is $120$, below the $138$ of our universe block (eigenvalues above $10^{-12}\lambda_{\max}$; \S\ref{sec:rank}).

\subsection{From App. A, appstatmethods.tex: the paragraph beginning ``The bootstrap varies events through distinct''}
% From app_statmethods.tex:245 @d55f7aad6; qualifier detail or procedure (DECISION-20261120a, second pass)
The bootstrap varies events through distinct data, signal-MC and background-MC sub-RNGs; universes vary model weights at a pinned GBDT seed; the training-seed block varies the GBDT seed over five settings, each on the same $20$ bootstrap replicas (at seed~42, replicas of the statistical block); the run-period block compares separate unfolds of three run periods. The bootstrap and the universes use the central value's LightGBM seed ($42$), and the training-seed settings include it.

\subsection{From App. A, appstatmethods.tex: the paragraph beginning ``The combined sqrtTrmathbfC of Tablereftab:uq is''}
% From app_statmethods.tex:262 @d55f7aad6; qualifier detail or procedure (DECISION-20261120a, second pass)
In the note, this text follows: ``The relative entries divide by the quoted central value, the LightGBM seed-42 central value of the universe sweep.''\par
Dividing the combined covariance by the exact-split cross-check instead gives a combined median of \SI{6.468}{\percent}. The universe, statistical and training-seed blocks are LightGBM ensembles, the first two at the central value's training seed and the third over five settings that include it (\S\ref{sec:method-config}). The universe and statistical blocks together give $5.976\%$, again evaluated on their sum rather than combined from the printed medians.

\subsection{From App. A, appstatmethods.tex: the paragraph beginning ``The paper reports mathrmd2sigma/(mathrmdpT,mathrmdpPar) on 14''}
% From app_statmethods.tex:282 @d55f7aad6; qualifier detail or procedure (DECISION-20261120a, second pass)
In the note, this text follows: ``Of these, $205$ have $x^{\paper}>0$ in \texttt{data\_result\_ptpl\_2D\_minerva\_inclusive\_6GeV.txt}.''\par
The paper suppresses the remaining $19$ through an efficiency cut. Every $\chi^{2}$ and pull excludes them on both sides.

\subsection{From App. A, appstatmethods.tex: the paragraph beginning ``A secondary strict-interior mask, used only''}
% From app_statmethods.tex:284 @d55f7aad6; qualifier detail or procedure (DECISION-20261120a, second pass)
A secondary strict-interior mask, used only as a sensitivity check, keeps cells with $\pT^{\text{high}} / \pPar^{\text{low}} \le \tan 20^{\circ}$, removing the $20$ reported cells that straddle the $\theta_{\mu}<20^{\circ}$ kinematic cut. The $185$-bin check tests whether the main tension comes from the boundary.
% carried with the text from app_statmethods.tex:285 @d55f7aad6
% Source: pre-freeze 2D_OMNIFOLD_RUN_LOG_ARCHIVE.md:50-59 (Phase 4, unreported bins), :716-726.

\subsection{From App. A, appstatmethods.tex: the paragraph beginning ``With boldsymbolDeltaequivmathbfxours-mathbfxpaper on the 205-bin set''}
% From app_statmethods.tex:305 @d55f7aad6; qualifier detail or procedure (DECISION-20261120a, second pass)
The comparison applies the pseudo-inverse of \S\ref{sec:invert} to the $205\times205$ block and computes the quadratic form. Options add OmniFold covariances via Eq.~\eqref{eq:blocksum} or subtract the paper's stat block to avoid counting it alongside the bootstrap (\S\ref{sec:rank}).

\subsection{From App. A, appstatmethods.tex: the paragraph beginning ``Setting mathrmndf=ntextreported requires nearly that many''}
% From app_statmethods.tex:309 @d55f7aad6; qualifier detail or procedure (DECISION-20261120a, second pass)
The choice rests on the $205$-bin 2D truncation scan rather than being assumed: the scan rises smoothly with the retained rank $r$, $0.39\,(r{=}50)\to1.04\,(100)\to1.76\,(180)\to2.27\,(205)$ (exact-split cross-check: $0.69\to2.35\to3.30\to3.66$), with no plateau or late jump, indicating an effective rank comparable to the dimension. This evidence is specific to 2D and does not transfer to N-D covariances (\S\ref{sec:nd-declarations}).

\subsection{From App. A, appstatmethods.tex: the paragraph beginning ``emph(i) Missing cross-covariance prevents calibration. For''}
% From app_statmethods.tex:338 @d55f7aad6; qualifier detail or procedure (DECISION-20261120a, second pass)
In the note, this text follows: ``and the published IBU-era and OmniFold results share MINERvA data and flux, detector and GENIE systematics.''\par
Without $C_{xy}$, the summed $C_{x}+C_{y}$ double-counts shared systematics and omits their correlation. Its net error has unknown size and sign.

\subsection{From App. A, appstatmethods.tex: the paragraph beginning ``emph(ii) Near-one realizations are not evidence''}
% From app_statmethods.tex:340 @d55f7aad6; qualifier detail or procedure (DECISION-20261120a, second pass)
Under calibrated assumptions, $\chi^{2}/\mathrm{ndf}$ has expectation $1$ and spread of order $\sqrt{2/\mathrm{ndf}}$, about $0.10$ for $\mathrm{ndf}=205$.

\subsection{From App. A, appstatmethods.tex: the paragraph beginning ``emph(iii) It does not address coverage''}
% From app_statmethods.tex:342 @d55f7aad6; qualifier detail or procedure (DECISION-20261120a, second pass)
In the note, this text follows: ``\emph{(iii) It does not address coverage} (\S\ref{sec:uq-2d-scope}, \S\ref{sec:coverage}).''\par
The covariance and the central value share the LightGBM estimator and, in the universe and statistical blocks, its training seed (\S\ref{sec:method-config}). Agreement claims use per-bin standardized differences and ratios (\S\ref{sec:pull}); this scalar is only a distance and ratio to the paper. The log-normal value $\chiCombinedLog$ depends on the definition given with the worked example.

\subsection{From App. A, appstatmethods.tex: the paragraph beginning ``The quoted chiCombined inverts the emphsum''}
% From app_statmethods.tex:349 @d55f7aad6; qualifier detail or procedure (DECISION-20261120a, second pass)
The quoted $\chiCombined$ inverts the \emph{sum} of the paper's covariance and ours, not our blocks alone. That sum contains the paper's \texttt{StatOnlyCov}, its published statistical-only covariance, which has rank $205$ and is near-diagonal (\S\ref{sec:rank}), so it places a variance floor under \emph{every} eigendirection of the combined matrix. It is the larger of the two statistical contributions. Our statistical block is $0.407$ of it in \emph{square-root} trace, and the square of that factor in \emph{trace} (operands in \S\ref{sec:rank}). Each ratio states which power of the trace it uses, because the two ratios describe the same comparison and differ by more than a factor of two. An external, published block, rather than any ensemble of ours, therefore sets the small-eigenvalue directions on which a finite-$N$ debiasing factor would act. This bounds the quoted number's exposure to declaration~(v) of \S\ref{sec:nd-declarations}.

\subsection{From App. A, appstatmethods.tex: the paragraph beginning ``Trace share alone cannot justify this''}
% From app_statmethods.tex:354 @d55f7aad6; qualifier detail or procedure (DECISION-20261120a, second pass)
Trace share alone cannot justify this protection: trace weights eigenvalues by $\lambda$, precision by $1/\lambda$, so even a negligible-trace block can dominate an inverse. Without the external floor, our finite-$N$ blocks fill $\mathbf{C}^{\syst}$'s null space (rank $138\to205$, eigenvalues above $10^{-12}\lambda_{\max}$) and, inverted directly, yield $\chi^{2}/\mathrm{ndf}=22.38$ (\S\ref{sec:rank}). Protection depends on the source of the smallest retained eigenvalues.

\subsection{From App. A, appstatmethods.tex: the paragraph beginning ``The paper covariance alone gives mean''}
% From app_statmethods.tex:382 @d55f7aad6; qualifier detail or procedure (DECISION-20261120a, second pass)
In the note, this text follows: ``The mean is a small difference between two estimates of the same data. It is not a finding of no bias.''\par
The standard deviation is slightly larger before our uncertainties are added.

\subsection{From App. A, appstatmethods.tex: the paragraph beginning ``This tests whether the emphshape of''}
% From app_statmethods.tex:387 @d55f7aad6; qualifier detail or procedure (DECISION-20261120a, second pass)
This tests whether the \emph{shape} of the $(\pT,\pPar)$ cross section matches the paper independently of absolute normalization, which depends on the flux integral, target nucleon count, and POT (each can drift by $\mathcal{O}(\%)$); the shape is a pure unfolding diagnostic.

\subsection{From App. A, appstatmethods.tex: the paragraph beginning ``A data-to-model chi2/mathrmndf with this analysis's''}
% From app_statmethods.tex:409 @d55f7aad6; qualifier detail or procedure (DECISION-20261120a, second pass)
In the note, this text follows: ``Every rank and condition number of our blocks quoted below uses this definition.''\par
Band count, not universe count, limits the rank: Eq.~\eqref{eq:matcov} gives each $\pm1\sigma$ pair rank $1$ and the $100$-universe PPFX band rank $\le 99$. PPFX is effectively rank-$1$ (\S\ref{sec:uq-universes}). The $43$ bands cannot span $205$ bins. The deficit does not come from missing bands.
% carried with the text from app_statmethods.tex:410 @d55f7aad6
% Source of rank and cond: budget_Q ranks (review NOTE 1). The ours-only script prints its own cond for the

\subsection{From App. A, appstatmethods.tex: the paragraph beginning ``A data-to-model chi2/mathrmndf with this analysis's''}
% From app_statmethods.tex:409 @d55f7aad6; qualifier detail or procedure (DECISION-20261120a, second pass)
In the note, this text follows: ``The paper's systematic-only covariance ($\texttt{TotalCov}-\texttt{StatOnlyCov}$) has rank $120$, below our $138$.''\par
We use the same \texttt{GetStandardSystematics} and propagate at least as many independent systematic directions.
% carried with the text from app_statmethods.tex:410 @d55f7aad6
% Source of rank and cond: budget_Q ranks (review NOTE 1). The ours-only script prints its own cond for the

\subsection{From App. A, appstatmethods.tex: the paragraph beginning ``Even at full rank (205) a''}
% From app_statmethods.tex:421 @d55f7aad6; qualifier detail or procedure (DECISION-20261120a, second pass)
The bootstrap block is smaller in \emph{global scale}: its square-root trace is $0.407$ of the paper's ($\sqrt{\Tr\mathbf{C}^{\boot}}=1.886\times10^{-40}$ vs $\sqrt{\Tr\,\texttt{StatOnlyCov}}=4.639\times10^{-40}$).

\subsection{From App. A, appstatmethods.tex: the paragraph beginning ``Even at full rank (205) a''}
% From app_statmethods.tex:421 @d55f7aad6; qualifier detail or procedure (DECISION-20261120a, second pass)
In the note, this text follows: ``The direct inverse of $\mathbf{C}^{\syst}+\mathbf{C}^{\boot}$ gives $\chi^{2}/\mathrm{ndf}=22.38$ ($21.96$ for all four blocks) for the LightGBM central value, against the paper covariance's $\chiPaper$ with its larger stat floor, so a usable $\chi^{2}_{\ours}$ depends on the \emph{statistical} block, not on more systematic bands.''\par
The gap is \emph{not} a missing MC-statistical term: our bootstrap already draws independent Poisson($1$) multipliers on data, signal MC and background MC per replica (\S\ref{sec:bootstrap}; Practical Guide 2507.09582 \S5.1). Three candidates remain. The first is a real difference in statistical efficiency between the unbinned reweighting and the published binned D'Agostini unfold. The second is a difference in what the paper's \texttt{StatOnlyCov} encompasses. The third is a regularization effect in our own estimator that suppresses dispersion without improving information. The split-bootstrap diagnostic used on the earlier rebuilt bootstrap has not been repeated for the quoted statistical block.

\subsection{From App. A, appstatmethods.tex: the paragraph beginning ``The data/MC split bootstrap was run''}
% From app_statmethods.tex:431 @d55f7aad6; qualifier detail or procedure (DECISION-20261120a, second pass)
In the note, this text follows: ``For the quoted statistical block, $\sqrt{\Tr\mathbf{C}^{\boot}}/\sqrt{\Tr\,\texttt{StatOnly}}=0.407$ and the per-bin median ratio of the two $\sigma$ is $0.425$.''\par
The fixed-truth test (\S\ref{sec:coverage}, \S\ref{sec:uq-boot}) does not decide between the explanations above.

\subsection{From App. A, appstatmethods.tex: the paragraph beginning ``For the exact-split cross-check and an''}
% From app_statmethods.tex:440 @d55f7aad6; qualifier detail or procedure (DECISION-20261120a, second pass)
Each method spans roughly two decades of $\chi^{2}/205$ across a two-decade scan of its regularization parameter, with no defensible single cutoff. The useful conclusion is diagnostic, not a number: the naive inverse is dominated by tiny, poorly-determined modes, and in the retained high-variance modes the central-value difference is not large. We therefore report direct per-bin pulls (\S\ref{sec:pull}) and this scan rather than a single unqualified $\chi^{2}_{\ours}$.

\subsection{From App. A, appstatmethods.tex: the paragraph beginning ``The combined distance (Srefsec:chisq) adds mathbfCboot''}
% From app_statmethods.tex:442 @d55f7aad6; qualifier detail or procedure (DECISION-20261120a, second pass)
In the note, this text follows: ``Subtracting that block uses $\texttt{TotalCov}-\texttt{StatOnlyCov}$ as baseline before our additions, changing $\chi^{2}/\mathrm{ndf}$ $\chiCombined\to \chiCombinedSubStat$.''\par
It over-corrects: our $\sqrt{\Tr\mathbf{C}^{\boot}}=1.886\times10^{-40}$ is $0.407$ of the paper's $\sqrt{\Tr\,\texttt{StatOnlyCov}}=4.639\times10^{-40}$, leaving no comparable replacement.

\subsection{From App. A, appstatmethods.tex: the paragraph beginning ``The known band-independence exception, the Bashyal''}
% From app_statmethods.tex:446 @d55f7aad6; qualifier detail or procedure (DECISION-20261120a, second pass)
In the note, this text follows: ``The quoted 2D covariance therefore remains the block sum, with the cross term as a diagnostic.''\par
A fitted $(\Phi_{\text{norm}},\Phi_{\text{shape}},\text{MuonE}_{\text{scale}})$ posterior from MAT-MINERvA could supply a paper-independent $\rho$. The collaboration uses the same truncated-spectral pseudo-inverse as $\chi^{2}_{\ours}$ (private communication), so the convention is reported to be collaboration practice rather than one introduced by this analysis. No published reference for it has been identified.

\subsection{From App. A, appstatmethods.tex: the paragraph beginning ``The combined chi2/mathrmndf=chiCombined is the descriptive''}
% From app_statmethods.tex:451 @d55f7aad6; qualifier detail or procedure (DECISION-20261120a, second pass)
The per-band sum is MAT-conformant, with the Flux band weighted as MAT weights it (\S\ref{sec:minerva}) and the $1.4\%$ normalization as a rank-1 band.

\subsection{From App. A, appstatmethods.tex: the paragraph beginning ``The variance methods of Srefsec:variance quantify''}
% From app_statmethods.tex:459 @d55f7aad6; qualifier detail or procedure (DECISION-20261120a, second pass)
The variance methods of \S\ref{sec:variance} quantify the scatter of the unfolded result. Closure tests measure its bias, a consistent offset from a known injected truth, which the variance methods cannot show.

\subsection{From App. A, appstatmethods.tex: the paragraph beginning ``OmniFold sees reco and truth (pT,pPar)''}
% From app_statmethods.tex:475 @d55f7aad6; qualifier detail or procedure (DECISION-20261120a, second pass)
In the note, this text follows: ``For median per-bin signed residual $\bar r(A)$, a linear fit gives $\bar r(A) \approx \kappa_{\text{leak}} A$, with $\kappa_{\text{leak}} = \partial \bar r/\partial A$, the fraction of hidden $\dpT$ bias entering the visible answer.''\par
The $3.0\%$ acceptance bound is looser than the truth-reweight bound because the weak $\dpT$--$\pT$ correlation prevents complete removal of the leakage.

\subsection{From App. A, appstatmethods.tex: the paragraph beginning ``Across 205 reported bins, combined OmniFold''}
% From app_statmethods.tex:491 @d55f7aad6; qualifier detail or procedure (DECISION-20261120a, second pass)
In the note, this text follows: ``Across $205$ reported bins, combined OmniFold $+$ paper uncertainties give $\chi^{2}/\mathrm{ndf} = \chiCombined$ and pull mean/standard deviation $0.034/0.389$. This distance is descriptive, without calibrated goodness-of-fit (\S\ref{sec:chisq}).''\par
Same-ensemble toy pulls are approximately Gaussian; the $0.048\%$ training-seed block is more than an order of magnitude below the systematic block.

\subsection{From App. A, appstatmethods.tex: the paragraph beginning ``Eq.eqrefeq:blocksum requires no cross-correlations between ML,''}
% From app_statmethods.tex:499 @d55f7aad6; qualifier detail or procedure (DECISION-20261120a, second pass)
In the note, this text follows: ``\textbf{Independence of error sources.} Eq.~\eqref{eq:blocksum} requires no cross-correlations between ML, statistical, systematic and run-period blocks.''\par
The bootstrap pins the GBDT seed, and universes pin the seed too. A common seed correlates the training noise and does not cancel it (\S\ref{sec:uq-2d-scope}); the training-seed block averages the same $20$ bootstrap replicas at each seed setting and includes the bootstrap\,$\times$\,seed interaction at its $1/20$ share. Among the first three, cross-terms are higher-order.

\subsection{From App. A, appstatmethods.tex: the paragraph beginning ``In reported bin i with bar''}
% From app_statmethods.tex:520 @d55f7aad6; qualifier detail or procedure (DECISION-20261120a, second pass)
In the note, this text follows: ``The fraction with $|r|\le1$ is \SI{\pullFrac}{\percent} (Gaussian reference $68.27\%$; $\Delta = +0.44\%$; median $68.50\%$), and $\langle|r|\rangle=\pullMeanAbs$ against $\sqrt{2/\pi}=\pullGauss$ (within $0.5\%$).''\par
The signed residual mean is $+0.006 \pm 0.082\,\sigma$, consistent with zero. The within-bin fraction is above \SI{65}{\percent} in \SI{97.6}{\percent} of bins (only $5$ of $205$ fall below). These show that the standardized toy residuals are approximately normal.

\subsection{From App. A, appstatmethods.tex: the paragraph beginning ``Because sigmai is estimated from the''}
% From app_statmethods.tex:522 @d55f7aad6; qualifier detail or procedure (DECISION-20261120a, second pass)
In the note, this text follows: ``Because $\sigma_i$ is estimated from the same toys that are evaluated, these figures diagnose standardized pulls and Gaussianity only and are not a coverage test of the bootstrap band.''\par
No coverage number is quoted from this ensemble, and the 2D coverage statements are the fixed-truth test of \S\ref{sec:coverage} and the held-out-tune test of \S\ref{sec:coverage-tunes}.
% carried with the text from app_statmethods.tex:523 @d55f7aad6
% Search: coverage-toys; coverage_summary; same-ensemble pull @ ddf16f10

\subsection{From App. A, appstatmethodsnd.tex: the paragraph beginning ``This appendix gives the conventions for''}
% From app_statmethods_nd.tex:5 @d55f7aad6; qualifier detail or procedure (DECISION-20261120a, second pass)
This appendix gives the conventions for covariances and comparisons beyond the 2D result: rank, ndf and finite-ensemble declarations for N-D covariances, the scalar five-dimensional diagnostic reductions and interval definitions, the extended-fiducial band and pull fractions, and inference from unfolded events and selected projections.



\subsection{From App. A, appstatmethodsnd.tex: the paragraph beginning ``The 2D justification in Srefsec:chisq does''}
% From app_statmethods_nd.tex:10 @d55f7aad6; qualifier detail or procedure (DECISION-20261120a, second pass)
In the note, this text follows: ``An earlier 5D covariance from the standard production chain has $10\,694$ reported bins, effective positive rank $\mathbf{263}$ and infinite numerical condition number. It is distinct from $C_Z$, the five-axis covariance of the \resFiveAxis{} (Table~\ref{tab:config-key}), and was built before $C_Z$ was chosen for use. Its $4825$-bin 5D$\to$4D projection retains rank $263$ exactly.''\par
The covariance is a sum of about 45 rank-one outer products, one for each two-endpoint MAT systematic band, and of the statistical and machine-learning (ML) sample blocks. This construction cannot span $10\,694$ directions and is rank deficient. Quoting $\chi^{2}/n_{\text{reported}}$ would understate tension by roughly $40\times$ in 5D or $18\times$ in 4D.

\subsection{From App. A, appstatmethodsnd.tex: the paragraph beginning ``For an N-D covariance, mathrmndf is''}
% From app_statmethods_nd.tex:12 @d55f7aad6; qualifier detail or procedure (DECISION-20261120a, second pass)
In the note, this text follows: ``For an N-D covariance, $\mathrm{ndf}$ is the retained rank.''\par
The reported-bin count $n_{\text{reported}}$ does not set it.

\subsection{From App. A, appstatmethodsnd.tex: the paragraph beginning ``For an N-D covariance, mathrmndf is''}
% From app_statmethods_nd.tex:12 @d55f7aad6; qualifier detail or procedure (DECISION-20261120a, second pass)
In the paragraph as the note printed it, this text followed: ``Where an N-D $\chi^{2}$ is quoted, it must state five things. They are (i) the inverse actually used, including the pseudo-inverse cutoff \texttt{rcond} (singular values below this fraction of the largest are discarded) or the truncation rank, (ii) the retained rank, as $\mathrm{ndf}$, and (iii) the rank-truncation scan of that covariance.''\par
That scan is expected to lack the smoothness that supports the 2D choice in \S\ref{sec:chisq}.

\subsection{From App. A, appstatmethodsnd.tex: the paragraph beginning ``Item (v) records the ensemble size''}
% From app_statmethods_nd.tex:16 @d55f7aad6; qualifier detail or procedure (DECISION-20261120a, second pass)
Item (v) records the ensemble size of each block. Without it, an uncorrected finite-ensemble bias could enter unnoticed.

\subsection{From App. A, appstatmethodsnd.tex: the paragraph beginning ``Item (v) records the ensemble size''}
% From app_statmethods_nd.tex:16 @d55f7aad6; qualifier detail or procedure (DECISION-20261120a, second pass)
In the note, this text follows: ``For this analysis's production convention, universe-mean centered with biased $1/N$ normalization, it is $(N-p-2)/N$.''\par
For fixed residuals an uncorrected bias inflates $\chi^{2}$ and overstates tension. Confidence regions from the same over-tight covariance become over-optimistic.

\subsection{From App. A, appstatmethodsnd.tex: the paragraph beginning ``No generic correction is prescribed. A''}
% From app_statmethods_nd.tex:19 @d55f7aad6; qualifier detail or procedure (DECISION-20261120a, second pass)
In the note, this text follows: ``No generic correction is prescribed.''\par
A sum of deterministic rank-one band terms and statistical and ML sample blocks, such as the earlier 5D covariance from the standard production chain, has no single debiasing factor.

\subsection{From App. A, appstatmethodsnd.tex: the paragraph beginning ``No generic correction is prescribed. A''}
% From app_statmethods_nd.tex:19 @d55f7aad6; qualifier detail or procedure (DECISION-20261120a, second pass)
In the note, this text follows: ``The standard factor also assumes independent Gaussian realizations and a data-independent truncation dimension, neither established for a data-dependent cut.''\par
The appropriate treatment depends on the covariance and inverse actually used, and is to be determined from the quantities recorded under item (v).

\subsection{From App. A, appstatmethodsnd.tex: the paragraph beginning ``The synthesis in Srefsec:gbdt-model-dependence reuses saved''}
% From app_statmethods_nd.tex:25 @d55f7aad6; qualifier detail or procedure (DECISION-20261120a, second pass)
The synthesis in \S\ref{sec:gbdt-model-dependence} reuses saved outputs. It involves no new unfolding, training or pseudo-experiments.

\subsection{From App. A, appstatmethodsnd.tex: the paragraph beginning ``The synthesis in Srefsec:gbdt-model-dependence reuses saved''}
% From app_statmethods_nd.tex:25 @d55f7aad6; qualifier detail or procedure (DECISION-20261120a, second pass)
In the note, this text follows: ``The saved trace functionals on the EW cells are differential in $(\Eavail,W)$ and are multiplied by the cell areas. Those on the J and H2 cells are already integrals.''\par
Direct tensor-slice integration checks the index-map sums and full-domain rate. Reporting boundaries coincide with fine edges.

\subsection{From App. A, appstatmethodsnd.tex: the paragraph beginning ``For experiment i and cell f,''}
% From app_statmethods_nd.tex:33 @d55f7aad6; qualifier detail or procedure (DECISION-20261120a, second pass)
In the note, this text follows: ``SD intervals use $s_f\sqrt{(n-1)/\chi^2_{0.975,n-1}}$ and $s_f\sqrt{(n-1)/\chi^2_{0.025,n-1}}$.''\par
These normal-model approximations condition on finite simulation and make no simultaneous claim across cells. A deterministic noise-free trace from one simulation receives no error band. Its cell-to-cell variation is not sampling uncertainty.

\subsection{From App. A, appstatmethodsnd.tex: the paragraph beginning ``The development assessment includes an interval''}
% From app_statmethods_nd.tex:35 @d55f7aad6; qualifier detail or procedure (DECISION-20261120a, second pass)
In the note, this text follows: ``The table gives integer hits, both empirical fractions, exact pointwise 95\% binomial intervals and fixed half-width relative to mean truth.''\par
It conditions on estimated bootstrap widths without propagating finite-replica uncertainty.

\subsection{From App. A, appstatmethodsnd.tex: the paragraph beginning ``The development assessment includes an interval''}
% From app_statmethods_nd.tex:35 @d55f7aad6; qualifier detail or procedure (DECISION-20261120a, second pass)
In the note, this text follows: ``Coverage fractions are not pooled across cells to certify them.''\par
The examples shown are a populated low-energy cell (EW7) and the corner cell EW41, already known to fail. They are not freshly selected validation targets.

\subsection{From App. A, appstatmethodsnd.tex: the paragraph beginning ``The rounding study keeps the fine-grid''}
% From app_statmethods_nd.tex:43 @d55f7aad6; qualifier detail or procedure (DECISION-20261120a, second pass)
In the note, this text follows: ``The estimated numerical share is $s_{\rm pair}^2/s_{\rm boot}^2$.''\par
The paired-model $\sqrt2$ factor introduces no physical noise source.

\subsection{From App. A, appstatmethodsnd.tex: the paragraph beginning ``For prior vertex v, data sensitivity''}
% From app_statmethods_nd.tex:47 @d55f7aad6; qualifier detail or procedure (DECISION-20261120a, second pass)
In the note, this text follows: ``For prior vertex $v$, data sensitivity is $d_{vf}=\hat f_{{\rm data},v}/\hat f_{{\rm data},0}-1$, with envelope $h_f=\max_v|d_{vf}|$.''\par
Its denominator differs from the simulated residual's.

\subsection{From App. A, appstatmethodsnd.tex: the paragraph beginning ``For prior vertex v, data sensitivity''}
% From app_statmethods_nd.tex:47 @d55f7aad6; qualifier detail or procedure (DECISION-20261120a, second pass)
In the note, this text follows: ``Neither $h_f$ nor an outer product of closure residuals is given a probability model.''\par
The reductions keep the per-cell values, the identity of each input and the exclusions, and distinguish complete outputs from checkpoints that carry no metadata. From a checkpoint, only the iterations that reproduce the originally recorded values are used. Later entries are excluded. The analyst cross-checked the reductions numerically.

\subsection{From App. A, appstatmethodsnd.tex: the paragraph beginning ``Unfolding correlates weighted events. Treating them''}
% From app_statmethods_nd.tex:58 @d55f7aad6; qualifier detail or procedure (DECISION-20261120a, second pass)
In the note, this text follows: ``Treating them as independent in a downstream likelihood can underestimate parameter uncertainties~\cite{Desai:2025correlated}.''\par
Nominal weights do not specify uncertainty. Coherent unfolding replicas could propagate it through the fit. Resampling unfolded rows independently is not equivalent to resampling inputs and repeating the learned estimator. Classifier and permutation goodness-of-fit tests also require null calibration for their fitting and data reuse. These are requirements on such inference, with no new validation result here.

\subsection{From App. A, appstatmethodsnd.tex: the paragraph beginning ``The uncertainty target matters: separate bin''}
% From app_statmethods_nd.tex:60 @d55f7aad6; qualifier detail or procedure (DECISION-20261120a, second pass)
In the note, this text follows: ``A fixed linear projection $y=P x$ propagates covariance as $C_y=P C_x P^{\top}$, with integrated bin widths in $P$ for densities.''\par
The identity neither calibrates that covariance nor accounts for regions chosen after inspecting data.

\subsection{From App. A, appstatmethodsnd.tex: the paragraph beginning ``Recent work constructs confidence intervals for''}
% From app_statmethods_nd.tex:62 @d55f7aad6; qualifier detail or procedure (DECISION-20261120a, second pass)
Recent work constructs confidence intervals for selected functionals in constrained inverse problems~\cite{Stanley:2025functionals} and joint confidence regions for multiple functionals~\cite{Batlle:2025simultaneous}, which could provide independent benchmarks for specified integrals or contrasts. Their guarantees depend on the forward model, noise law, and parameter constraints and do not transfer automatically to a learned response with finite simulation and detector nuisances. A rank-deficient forward operator is also distinct from a rank-deficient sample covariance. Neither truncating the latter nor validating a few projections certifies the full distribution. Such a benchmark would supplement the scalar five-dimensional uncertainty construction that a measurement would require. It would not replace that construction.

\subsection{From App. A, appstatmethodsnd.tex: the paragraph beginning ``Each calibration pseudo-experiment splits the simulation''}
% From app_statmethods_nd.tex:69 @d55f7aad6; qualifier detail or procedure (DECISION-20261120a, second pass)
In the note, this text follows: ``\paragraph{Generative model of one pseudo-experiment.} Each calibration pseudo-experiment splits the simulation at random into two halves (Table~\ref{tab:joint5d-genmodel}) and refits the background refinement. Its truth equals $G$ on every fine cell that the source half populates.''\par
Cells it leaves empty lower a coarse cell by at most \SI{0.2}{\percent}, or about \SI{4}{\percent} of the null standard deviation (the per-experiment spread is below \SI{1.5}{\percent} of it).

\subsection{From App. A, appstatmethodsnd.tex: the paragraph beginning ``The data's rounding noise, and the''}
% From app_statmethods_nd.tex:71 @d55f7aad6; qualifier detail or procedure (DECISION-20261120a, second pass)
In the note, this text follows: ``The prediction's variance in that draw is its own per-cell sum of squared weights. That diagonal model is an assumption of the metric.''\par
Disjoint cells do not by themselves make a prediction's Monte Carlo errors independent. The background source and its template are Poisson draws from disjoint halves of the background sample. The detector weight bands include the GEANT4 hadron-interaction reweights.

\subsection{From App. A, appstatmethodsnd.tex: the paragraph beginning ``emphProbability model. Fix one test, and''}
% From app_statmethods_nd.tex:132 @d55f7aad6; qualifier detail or procedure (DECISION-20261120a, second pass)
In the note, this text follows: ``Assumption A2 is that, given the data, these are independent Bernoulli draws with a common probability $p$, the calibrated tail probability of the observed statistic.''\par
This requires the lost pseudo-experiments to be lost independently of their outcomes.

\subsection{From App. A, appstatmethodsnd.tex: the paragraph beginning ``emphProbability model. Fix one test, and''}
% From app_statmethods_nd.tex:132 @d55f7aad6; qualifier detail or procedure (DECISION-20261120a, second pass)
In the note, this text follows: ``Assumption A1 is that, for a true null, $p$ is super-uniform over repetitions of the experiment, $P(p\le u)\le u$, with the nuisance parameters drawn from their priors.''\par
The size statement is therefore averaged over the nuisance priors.

\subsection{From App. A, appstatmethodsnd.tex: the paragraph beginning ``emphHolm reduction without independence. Let m0''}
% From app_statmethods_nd.tex:140 @d55f7aad6; qualifier detail or procedure (DECISION-20261120a, second pass)
In the note, this text follows: ``A familywise error therefore requires some true null to cross below $\alpha/m_{0}$, and by the union bound $\mathrm{FWER}\le m_{0}\,R(\alpha/m_{0})\le\max_{j\le m}j\,R(\alpha/j)$.''\par
The argument uses neither independence between tests nor the order of the Holm steps, so the shared ensemble of a null's total and shape tests, their unequal stopping counts, and the ranking by $(k+1)/(B+1)$ with its ties do not affect it.

\subsection{From App. A, appstatmethodsnd.tex: the paragraph beginning ``The largest ratio R(c)/c, at j=1,''}
% From app_statmethods_nd.tex:142 @d55f7aad6; qualifier detail or procedure (DECISION-20261120a, second pass)
In the note, this text follows: ``The largest ratio $R(c)/c$, at $j=1$, lies in $[0.8892,0.8902]$.''\par
It decreases to $[0.5138,0.5148]$ at $j=10$.

\subsection{From Part 2, secmethod.tex: the paragraph beginning ``paragraphTraining details. Each classifier is trained''}
% From sec_method.tex:17 @fdd67e926; training detail (DECISION-20261120a records route; the director 2026-11-20)
The two classes keep their physical weights, with no class balancing and no renormalization between iterations. The learned ratio $p/(1-p)$ therefore carries the normalization, with $p$ restricted to $[10^{-6},1-10^{-6}]$.

\subsection{From Part 2, secmethod.tex: the paragraph beginning ``paragraphTraining details. Each classifier is trained''}
% From sec_method.tex:17 @fdd67e926; training detail (DECISION-20261120a records route; the director 2026-11-20)
In the note, this text follows: ``In the refined configurations, a classifier separating positive from negative weights, trained with $|w|$, gives each event the weight $|w|\max(0,2g-1)$.''\par
For the \estMultiD{}, the per-event Kish effective sample size and weight tails have been computed on the sixteen reported cells~\cite{UnfoldedWeightsESS2026}. They describe the unfolded simulated-event weights and are not an uncertainty, coverage or validity statement.
